\section{Introduction}
\label{sec:introduction}

Brain tumors represent one of the most critical neurological disorders,
requiring accurate and timely diagnosis to improve treatment planning and
patient outcomes. Magnetic Resonance Imaging (MRI) is widely used for brain
tumor assessment because of its superior soft-tissue contrast and
non-invasive nature. However, manual interpretation of MRI scans remains
challenging due to the substantial variability in tumor appearance, shape,
size, location, and texture. Consequently, computer-aided diagnosis (CAD)
systems have emerged as valuable tools for assisting radiologists in the
classification of brain tumors~\cite{ottoni2025}.

Recent advances in artificial intelligence have significantly improved
medical image analysis. Deep learning models, particularly Convolutional
Neural Networks (CNNs), have achieved remarkable performance in brain tumor
classification tasks~\cite{ottoni2025}. Nevertheless, these approaches often
require large annotated datasets, extensive computational resources, and
specialized hardware for training. Furthermore, their decision-making
process is frequently difficult to interpret, which may limit their adoption
in clinical environments where transparency and explainability are important
considerations~\cite{ottoni2025}.

In contrast, handcrafted feature-based approaches remain attractive for
small- and medium-sized medical imaging datasets. Texture and structural
descriptors such as the Gray-Level Co-occurrence Matrix
(GLCM)~\cite{haralick1973}, Local Binary Patterns (LBP)~\cite{ojala2002},
the Discrete Wavelet Transform (DWT)~\cite{mallat1989}, and the Histogram of
Oriented Gradients (HOG)~\cite{dalal2005} provide interpretable
representations of image characteristics while maintaining relatively low
computational requirements. Numerous studies have demonstrated that
carefully designed handcrafted features combined with classical machine
learning classifiers can achieve competitive performance in brain tumor
classification~\cite{pattanaik2022,dheepak2024,nawaz2022,basthikodi2024}.

Despite their effectiveness, most existing studies rely on default descriptor
parameters or optimize feature extractors using classifier-dependent
criteria~\cite{pattanaik2022,dheepak2024,basthikodi2024}. In such settings,
the quality of the extracted features becomes intertwined with the behavior
of a specific classifier, making it difficult to determine whether
performance improvements originate from the feature representation itself or
from the learning algorithm. Moreover, many studies evaluate only a limited
number of classifiers and rarely provide comprehensive statistical analyses
to verify the significance of observed performance
differences~\cite{pattanaik2022,ottoni2025}.

To address these limitations, this work proposes a two-phase framework for
brain tumor MRI classification. In the first phase, handcrafted feature
extractors are configured using a classifier-independent optimization
strategy based on feature separability. Three complementary separability
metrics---Fisher Discriminant Ratio (FDR), Mutual Information (MI), and
Davies--Bouldin Index (DBI)---are employed to evaluate and rank candidate
parameter configurations for GLCM, LBP, DWT, and HOG descriptors. This
approach enables the selection of feature representations that maximize class
discrimination before any classifier training is performed.

In the second phase, the optimized feature representations are used to
conduct a comprehensive benchmarking study of eight classical machine
learning classifiers, namely Support Vector Machine (SVM)~\cite{cortes1995},
K-Nearest Neighbors (KNN), Random Forest (RF)~\cite{breiman2001},
XGBoost~\cite{chen2016}, LightGBM~\cite{ke2017}, Logistic Regression (LR),
Multi-Layer Perceptron (MLP), and Extra Trees (ET). The benchmarking process
evaluates multiple feature combinations and employs a rigorous experimental
protocol based on stratified cross-validation and hyperparameter
optimization. In addition, statistical validation techniques are applied to
assess the reliability and significance of the obtained results.

The main contributions of this paper can be summarized as follows:
\begin{enumerate}
    \item A separability-guided methodology for classifier-independent
    optimization of handcrafted feature extractors in brain tumor MRI
    classification.
    \item A systematic evaluation of GLCM, LBP, DWT, and HOG descriptor
    configurations using Fisher Discriminant Ratio, Mutual Information, and
    Davies--Bouldin Index.
    \item A comprehensive benchmark of eight classical machine learning
    classifiers across multiple optimized feature representations.
    \item An analysis of feature-group importance that identifies which
    descriptor families contribute most to classification performance.
    % NOTE: softened from "the relationship between descriptor separability
    % and classification performance"; restore only if that analysis is added.
    \item A rigorous statistical validation framework incorporating Cohen's
    kappa coefficient~\cite{cohen1960}, bootstrap confidence intervals,
    McNemar's test~\cite{mcnemar1947}, the Wilcoxon signed-rank test, and
    Friedman--Nemenyi analysis~\cite{demsar2006}.
\end{enumerate}

The results demonstrate that separability-guided feature configuration
combined with statistically validated classifier benchmarking can provide
competitive performance for brain tumor MRI classification while maintaining
the interpretability and computational efficiency associated with handcrafted
feature engineering.
% NOTE: "highly competitive" -> "competitive" until final Phase 2 numbers are in.
